% Generated draft: requires mathematical and source review.
\documentclass[11pt]{article}
\usepackage[margin=25mm]{geometry}
\usepackage{fontspec}
\setmainfont{Latin Modern Roman}
\usepackage{amsmath,amssymb,mathtools,unicode-math}
\setmathfont{Latin Modern Math}
\usepackage{enumitem,needspace,xurl,hyperref}
\setlength{\parindent}{0pt}
\setlength{\parskip}{5pt}
\setlength{\emergencystretch}{3em}
\newcommand{\N}{\mathbb{N}}
\newcommand{\Z}{\mathbb{Z}}
\newcommand{\Q}{\mathbb{Q}}
\newcommand{\R}{\mathbb{R}}
\newcommand{\C}{\mathbb{C}}
\newcommand{\F}{\mathbb{F}}
\newcommand{\A}{\mathbb{A}}
\newcommand{\PP}{\mathbb{P}}
\newcommand{\E}{\mathbb{E}}
\newcommand{\eps}{\varepsilon}
\newcommand{\dd}{\,\mathrm{d}}
\begin{document}
\section*{UnsolvedMath ID: 30000883}
\textbf{Source problem code:} catalogue entry
\subsection{Definitions and mathematical statement}
For a group $G$, let $\Z G$ be its integral group ring and let $U=U(\Z G)$ be its full unit group, containing the canonical copy of $G$. Define the finite-conjugacy center and normalizer by \[\Delta(U)=\{u\in U:[U:C_U(u)]<\infty\},\qquad N_U(G)=\{u\in U:uGu^{-1}=G\}.\] Thus finite conjugacy concerns the size of a conjugacy class, not finite order of $u$. Define the upper central series by $Z_0(U)=1$, $Z_{\gamma+1}(U)/Z_\gamma(U)=Z(U/Z_\gamma(U))$, and union at limit ordinals. Its stable terminal term is the hypercenter, denoted $Z_\infty(U)$ in the source. The first two questions are whether \[Z_\infty(U)\subseteq\Delta(U),\qquad\Delta(U)\subseteq N_U(G).\] The extract does not state whether this first pair assumes finite $G$; that assumption is explicit only in the next pair.\par Now let $G$ be finite and define the normalized units by \[V(\Z G)=\ker\bigl(\varepsilon:U(\Z G)\to\{1,-1\}\bigr),\qquad\varepsilon\Bigl(\sum_g a_gg\Bigr)=\sum_g a_g.\] A section of a group is a quotient $B/N$ with $N\trianglelefteq B$ and $B$ a subgroup. Is every simple section of every finite subgroup $A\leq V(\Z G)$ isomorphic to a section of $G$? Explicitly, if $N\trianglelefteq B\leq A$ and $B/N$ is simple, must there exist $N'\trianglelefteq B'\leq G$ with $B/N\simeq B'/N'$? The source does not restrict simple to nonabelian simple groups. Finally, for any group $K$, define its prime graph $\Pi(K)$ to have as vertices the primes $p$ for which $K$ contains an element of order $p$, with distinct $p,q$ adjacent when $K$ contains an element of order $pq$. Does \[\Pi(V(\Z G))=\Pi(G)\] hold, including both vertex sets and edges? This definition uses finite-order elements even though the full unit group is infinite. The supplied 2026-08-21 literature triage labels the record open and reports many known families but no general answer to the prime-graph and subgroup-section problems.
\subsection{Short English statement}
Ask for two inclusions among the hypercenter, finite-conjugacy center, and normalizer in the full unit group. Separately, for finite $G$, ask whether simple sections of finite normalized-unit subgroups occur as sections of $G$ and whether their prime graphs coincide.
\subsection{Sources}
\begin{itemize}\item Local UnsolvedMath registry, permanent ID 30000883, source problem catalogue entry, \emph{Hypercenter and Prime-Graph Problems for Integral Group Rings}, proposed year 2007; no individual proposer supplied.\item Supplied citation: \emph{Mini-Workshop: Arithmetik von Gruppenringen} (2008), Oberwolfach Reports, DOI 10.4171/owr/2007/55, pp.\ 3209--3240: \url{https://doi.org/10.4171/owr/2007/55}. Supplied literature label: original Oberwolfach report.\item Supplied literature label: Gruenberg--Kegel graph reduction, \url{https://arxiv.org/abs/1612.05025}.\item Supplied literature label: prime graph status and results, \url{https://www.sciencedirect.com/science/article/pii/S0747717119300161}.\end{itemize}
\end{document}
